\DocumentMetadata{
  pdfversion=2.0,
  pdfstandard=ua-2,
  lang=en-US,
  tagging=on,
  tagging-setup={math/setup=mathml-SE,tabsorder=structure}
}
\documentclass[11pt,letterpaper]{article}
\usepackage[margin=0.68in,includefoot]{geometry}
\usepackage{newcomputermodern}
\usepackage{amsmath,amscd,mathtools}
\usepackage{tikz,adjustbox}
\providecommand{\square}{\mdlgwhtsquare}
\usepackage[hidelinks]{hyperref}
\hypersetup{
  pdftitle={Algebra PhD Exam, May 2011},
  pdfauthor={University of Florida Department of Mathematics},
  pdfsubject={Graduate mathematics examination},
  pdfdisplaydoctitle=true,
  bookmarksopen=true
}
\setcounter{secnumdepth}{0}
\setlength{\parindent}{0pt}
\setlength{\parskip}{0.28em}
\setlength{\emergencystretch}{3em}
\setlength{\leftmargini}{2.15em}
\setlength{\leftmarginii}{2.65em}
\setlength{\leftmarginiii}{3em}
\pagestyle{empty}
\raggedbottom
\allowdisplaybreaks
\NewTaggingSocketPlug{block/list/label}{examproblem}{%
  \tagstructbegin{tag=Lbl}\tagstructend%
  \tagstructbegin{tag=\LBody}%
  \tagstructbegin{tag=\ProblemHeadingTag,title-o={\ProblemHeadingText}}%
  \tagmcbegin{tag=\ProblemHeadingTag,actualtext={\ProblemHeadingText}}%
  #2%
  \tagmcend%
  \tagstructend%
}
\newenvironment{examproblems}[2]{%
  \def\ProblemHeadingTag{#2}%
  \AssignTaggingSocketPlug{block/list/label}{examproblem}%
  \begin{enumerate}%
  \setcounter{enumi}{#1}%
  \renewcommand{\labelenumi}{\textbf{\theenumi.}}%
  \setlength{\topsep}{0.35em}%
  \setlength{\partopsep}{0pt}%
  \setlength{\itemsep}{0.48em}%
  \setlength{\parsep}{0.18em}%
}{\end{enumerate}}
\newenvironment{examsubparts}{%
  \AssignTaggingSocketPlug{block/list/label}{default}%
  \begin{enumerate}%
  \setlength{\topsep}{0.18em}%
  \setlength{\partopsep}{0pt}%
  \setlength{\itemsep}{0.16em}%
  \setlength{\parsep}{0.1em}%
}{\end{enumerate}}
\newenvironment{examinstructions}{%
  \par\begingroup\itshape
}{%
  \par\endgroup\vspace{0.18em}
}
\makeatletter
\renewcommand\subsection{\@startsection{subsection}{2}{0pt}%
  {-1.25ex plus -.25ex minus -.1ex}%
  {0.55ex plus .1ex}%
  {\normalfont\large\bfseries}}
\renewcommand{\@maketitle}{%
  \newpage\null\vskip -1em%
  {\footnotesize\bfseries
    \href{https://www.ufl.edu/}{University of Florida}, \href{https://math.ufl.edu/}{Department of Mathematics}\par}%
  \vskip -0.5em%
  \tagstructbegin{tag=H1,title={Algebra PhD Exam, May 2011}}%
  \tagpdfparaOff%
  \tagmcbegin{tag=H1}%
    {\LARGE\bfseries\href{https://gma.math.ufl.edu/exams/algebra-phd/}{Algebra PhD Exam}, May 2011\par}%
  \tagmcend%
  \tagpdfparaOn%
  \tagstructend%
  \vskip 0.25em%
  \tagmcbegin{artifact}%
    \hrule height 1pt%
  \tagmcend%
  \vskip 1em%
}
\makeatother
\title{Algebra PhD Exam, May 2011}
\author{}
\date{}
\begin{document}
\maketitle
\thispagestyle{empty}
\begin{examinstructions}
Answer seven problems. Write your answers clearly in complete English sentences. You may quote results (within reason) as long as you state them clearly.
\end{examinstructions}
\begin{examproblems}{0}{H2}
\def\ProblemHeadingText{Problem 1.}
\item
Let~\(K\) be a finite field, and let~\(F\) be a subfield of~\(K\). Prove that the extension \(K/F\) is Galois.
\def\ProblemHeadingText{Problem 2.}
\item
This problem has three parts.
\begin{examsubparts}
\item[\textnormal{(a)}]
Determine a splitting field~\(K\) for the polynomial \(X^3-2\in\mathbb{Q}[X]\).
\item[\textnormal{(b)}]
Determine \(\operatorname{Gal}(K/\mathbb{Q})\) and give an explicit description of the action of \(\operatorname{Gal}(K/\mathbb{Q})\) on~\(K\).
\item[\textnormal{(c)}]
For each subgroup~\(H\) of \(\operatorname{Gal}(K/\mathbb{Q})\) find the subfield of~\(K\) fixed by~\(H\).
\end{examsubparts}
\def\ProblemHeadingText{Problem 3.}
\item
Let~\(A\) and~\(B\) be abelian groups, and \(T=A\otimes_{\mathbb{Z}}B\). Let \(\phi:A\to A\) be a group homomorphism. Prove that there is a group homomorphism \(\psi:T\to T\) such that, for all \(a\in A\) and \(b\in B\), we have \(\psi(a\otimes_{\mathbb{Z}}b)=\phi(a)\otimes_{\mathbb{Z}}b\).
\def\ProblemHeadingText{Problem 4.}
\item
This problem has two parts.
\begin{examsubparts}
\item[\textnormal{(a)}]
Let~\(\mathcal{C}\) be a concrete category and let~\(S\) be a set. Define what it means to say that an object in~\(\mathcal{C}\) is free for~\(S\).
\item[\textnormal{(b)}]
Prove that any two objects in~\(\mathcal{C}\) which are free for~\(S\) are equivalent in~\(\mathcal{C}\).
\end{examsubparts}
\def\ProblemHeadingText{Problem 5.}
\item
Let~\(R\) be a ring.
\begin{examsubparts}
\item[\textnormal{(a)}]
Define what it means for an~\(R\)-module~\(M\) to be projective.
\item[\textnormal{(b)}]
Prove that any free~\(R\)-module is projective.
\end{examsubparts}
\def\ProblemHeadingText{Problem 6.}
\item
Let~\(F\) be a free group on a set~\(S\). Prove that~\(F\) is abelian if and only if \(|S|\leq 1\).
\def\ProblemHeadingText{Problem 7.}
\item
State and prove Hilbert’s Basis Theorem.
\def\ProblemHeadingText{Problem 8.}
\item
Let~\(R\) be a commutative ring with identity. Let~\(S\) be a non-empty subset of~\(R\) which is closed under multiplication and is such that \(0\notin S\). Prove that there exists a prime ideal~\(P\) of~\(R\) such that \(P\cap S=\varnothing\).
\def\ProblemHeadingText{Problem 9.}
\item
Let~\(R\) be a Dedekind domain, and let~\(I\) be a non-zero ideal of~\(R\). Prove that \(R/I\) is Artinian.
\def\ProblemHeadingText{Problem 10.}
\item
This problem has three parts.
\begin{examsubparts}
\item[\textnormal{(a)}]
Let~\(R\) be an integral domain. Define what is a fractional ideal of~\(R\).
\item[\textnormal{(b)}]
Define what it means for a fractional ideal of~\(R\) to be invertible.
\item[\textnormal{(c)}]
Prove from your definition that if~\(R\) is a principal ideal domain then every fractional ideal of~\(R\) is invertible.
\end{examsubparts}
\def\ProblemHeadingText{Problem 11.}
\item
Let~\(R\) and~\(S\) be non-commutative semisimple rings with \(|R|=|S|=2^2\cdot 3^4\cdot 5\). Does it follow that~\(R\) is isomorphic to~\(S\)? Justify your answer.
\end{examproblems}
\end{document}
